\documentclass[10pt,a4paper]{amsart}
\usepackage[margin=19mm]{geometry}
\usepackage{amsmath,amssymb,mathtools}
\usepackage[scr=rsfs,cal=euler]{mathalfa}
\usepackage{xfrac}
\usepackage[unicode,hidelinks,bookmarks=false]{hyperref}
\newcommand{\ie}{i.e.\ }
\newcommand{\cf}{cf.\ }
\newcommand{\resp}{respectively\ }
\newcommand{\Subsection}{\subsection}
\providecommand{\nobreakdash}{\nobreak\hbox{-}}
\setlength{\emergencystretch}{2em}
\multlinegap0pt
\usepackage{bm}
\usepackage{bbm}
\usepackage{dsfont}
\usepackage{xspace}
\usepackage{cancel}
\usepackage{mathtools}
\usepackage{amsthm}
\makeatletter
\@ifundefined{theorem}{\newtheorem{theorem}{Theorem}}{}
\@ifundefined{corollary}{\newtheorem{corollary}[theorem]{Corollary}}{}
\makeatother
\def\NN{\mathbb N}
\def\aa{\alpha}
\def\aim{\text{targeting}{}}
\newcommand{\free}[1]{\mathcal{F}(#1)} 
\title[Hypercyclicity and Lipschitz-free operators]{Hypercyclicity and Lipschitz-free operators}
\author{Christian Cobollo, Romuald Ernst, Quentin Menet, Alfred Peris}
\date{Source: arXiv:2609.17874v1 (CC BY 4.0)}
\begin{document}
\maketitle

\noindent\emph{Source and licence.} Hypercyclicity and Lipschitz-free operators. Version arXiv:2609.17874v1 (CC BY 4.0), distributed under the Creative Commons Attribution 4.0 International licence, as recorded in the arXiv licence metadata for this exact version and shown on the abstract page https://arxiv.org/abs/2609.17874v1. Full rights notice: licenses/arxiv-2609.17874-CC-BY-4.0.txt.\par\smallskip

\noindent\textit{Editorial scope.} Counted unit: the Theorem whose source label is \texttt{hypdeltas} in the article (source0.tex lines 959--966, source label \texttt{hypdeltas}), delivered together with the article's own complete original proof of it (source0.tex lines 967--1047, one \texttt{proof} environment of 81 lines). No mathematics is written, repaired or re-derived: statement and proof are the article's own text line for line. Required context retained uncounted: thm:aim-WHC (lines 767--773); cor:delta-basis (lines 907--912). Every definition, display and label of those lines is retained unchanged. Omitted from this package and not redistributed: 1--766, 774--906, 913--958, 1048--1298. Nothing inside a retained excerpt is omitted or altered: every retained body is delivered exactly as the article's lines are written, so no omission receipt of any kind accompanies this package. Citations: the retained excerpts carry no citation command at all, so no bibliography entry of the article is transcribed and no reference is redistributed. Reference adaptation: none was needed and none was made; every reference inside the delivered unit resolves inside it, so no unresolved reference or citation is produced. Printed slips kept literal: no printed word, symbol, hypothesis or number of the counted statement or of its proof was altered. Layout and class: the article's own class is \texttt{amsart} with packages including geometry, amssymb, amsmath, amsthm, enumerate, inputenc, babel, xcolor, url, graphicx, tikz, csquotes, biblatex, hyperref; the wrapper is the lane's pinned \texttt{amsart} editorial wrapper at 10pt on A4, which restates the theorem-like environments the retained text needs and adds the article's own macro definitions that the retained text needs (\texttt{\textbackslash NN}, \texttt{\textbackslash aa}, \texttt{\textbackslash aim}, \texttt{\textbackslash free}), with extra wrapper packages (bm, bbm, dsfont, xspace, cancel, mathtools). The delivered numbering is the unit's own. Assets: no retained span contains a figure environment, a graphics-inclusion command, a PDF-page inclusion, a table or a TikZ picture; no image file is copied into this package, the package declares no asset dependency and no image file is redistributed with it. Licence: version arXiv:2609.17874v1 of this article is distributed under the Creative Commons Attribution 4.0 International licence, as recorded in the arXiv licence metadata for this exact version and shown on the abstract page https://arxiv.org/abs/2609.17874v1. A full rights notice is delivered at licenses/arxiv-2609.17874-CC-BY-4.0.txt. Characters: every retained excerpt is pure ASCII, so the delivered engine, XeLaTeX with its default Latin Modern text font, reports no missing character.
\begin{theorem}  \label{thm:aim-WHC}
Let $T\colon X\to X$ be a linearization on a topological vector space $X$ of a continuous map $f\colon Y\to Y$ on a uniform space $Y$ with fixed point $x_0$ such that the embedding $j\colon Y\to X$ is uniformly continuous. We have: 
\begin{enumerate}
    \item If $f$ has the targeting property around $x_0$ then $T$ is weakly mixing.
    \item If $f$ has the strict targeting property anon-empty open sets round $x_0$ th en $T$ is mixing.
\end{enumerate}
\end{theorem}

\begin{corollary}\label{cor:delta-basis}
Let $(M,d)$ be metric space with a distinguished point $0$. The space $\mathcal{F}(M)$ admits a Schauder basis given by deltas if and only if $M$ is at most countable and radially discrete with respect to $0$.
Moreover, such a basis is given by the deltas of any  enumeration of $M\backslash\{0\}$, and it is unconditional, since the norm on $\mathcal{F}(M)$ is equivalent to the norm 
\[\left\|\sum_{x\in M\backslash\{0\}}\aa_x \delta_{x}\right\|_1=\sum_{x\in M\backslash\{0\}}\aa_x d(x,0).\]
In particular, $\free{M}$ is isomorphic to $\ell^1(\NN)$.
\end{corollary}

\begin{theorem}\label{hypdeltas}
Let $(M,d)$ be a metric space with a distinguished point $0$ and $f\in \text{Lip}_0(M,M)$. If $M$ is countable and radially discrete with respect to $0$, then the following assertions are equivalent:
\begin{enumerate}
\item $T_f$ is hypercyclic on $\free{M}$;
\item $T_f$ is weakly mixing on $\free{M}$;
\item $f$ has the \aim{} property around $0\in M$.
\end{enumerate}
\end{theorem}

\begin{proof}
Let $(x_n)_{n\geq1}$ be an enumeration of $M\backslash\{0\}$. We know by Corollary \ref{cor:delta-basis} that $(\delta_{x_n})_{n\geq1}$ is a Schauder basis in $\free{M}$.
Moreover, we have (3) $\Rightarrow$ (2) $\Rightarrow$ (1) by Theorem~\ref{thm:aim-WHC}.\\


(1) $\Rightarrow$ (3). Assume that $T_f$ is hypercyclic and thus topologically transitive.\\


Let $\varepsilon>0$, $k\ge 1$ and $m_0\ge 1$. We want to show that there exists $m\ge m_0$ such that for every $1\le j\le k$, 
\begin{itemize}
\item there exists $x\in M$ such that $d(x,0)<\varepsilon$  and $d(f^m(x),f^m(x_j))<\varepsilon$,
\item there exists $x'\in M$ such that $d(x',0)<\varepsilon$ and $f^m(x')=x_j$.
\end{itemize}
This will imply that $f$ has the targeting property around $0$.\\

To this end, we let $M_k=\{x_1,\dots,x_k\}$ and
\[\mu_k=\sum_{j=1}^k 3^j\delta_{x_j} \quad \text{and}\quad \nu_k=\sum_{j=1}^k 3^{k+j}\delta_{x_j}.\]
Therefore, by topological transitivity of $T_f$, there exist an integer $m\ge m_0$ and $\rho\in \mathcal{F}(M)$ such that 
\[\|\rho-\mu_k\|_1<\gamma \quad \text{and}\quad \|T_f^m\rho-\nu_k\|_1<\gamma\]
where $\gamma$ is given by
\[\gamma=\min\left\{\frac{\min\{d(x_j,0):1\le j\le k\}}{k},\varepsilon\right\}.\]
By density of $\text{vect}\{\delta_{x_n}: n\ge 1\}$, we can assume that $\rho=\sum_{j=1}^N a_j \delta_{x_j}$ with $N\ge k$. We then have
\[\left\|\sum_{j=1}^N a_j \delta_{x_j}-\sum_{j=1}^k 3^j\delta_{x_j}\right\|_1<\gamma \quad \text{and}\quad  \left\|\sum_{j=1}^N a_j \delta_{f^m(x_j)}-\sum_{j=1}^k 3^{k+j}\delta_{x_j}\right\|_1<\gamma.\]

In particular, we deduce by using the value of $\gamma$ that
\begin{enumerate}
\item for every $1\le j\le k$, $\displaystyle |a_j-3^j|<\frac{1}{k}$,
\item  $\displaystyle \sum_{k< j\le N}|a_j|d(x_j,0)<\varepsilon$,
\item for every $1\le j\le k$, $\displaystyle\left|\sum_{1\le i\le N:f^m(x_i)=x_j}a_i-3^{k+j}\right|<1$,
\item for every $x\in M\backslash M_k$,
\[\displaystyle\left|\sum_{1\le i\le N:f^m(x_i)=x}a_i\right| d(x,0)<\varepsilon.\]
\end{enumerate}
It follows from $(1)$ that for every set $F\subset \{1,\dots,k\}$ and every $1\le j\le k$, we have
\[\left|\sum_{i\in F} a_i-3^{k+j}\right|> 2.\]

Indeed, by (1), we have
\begin{align*}
     \left|\sum_{i\in F}a_i\right| < \sum_{i=1}^k\left(3^i+\frac{1}{k}\right)= \dfrac{3^{k+1}-1}{2},
\end{align*}
and from this follows
\begin{align*}
    \left|\sum_{i\in F} a_i-3^{k+j}\right | \geq 3^{k+j}-\dfrac{3^{k+1}-1}{2} >\frac{3^{k+1}}{2}> 2.
\end{align*}
We then deduce from $(1)$ and $(3)$ that for every $1\le j\le k$, 
\[\sum_{i\in ]k,N]:f^m(x_i)=x_j}|a_i|\ge \left|\sum_{i\in ]k,N]:f^m(x_i)=x_j}a_i\right|>1\]

since
\begin{align*}
\left|\sum_{i\in ]k,N]:f^m(x_i)=x_j}a_i\right|&=\left|\left(\sum_{1\le i\le N:f^m(x_i)=x_j}a_i-3^{k+j}\right)-\left(\sum_{1\le i\le k:f^m(x_i)=x_j}a_i-3^{k+j}\right) \right|\\
&\ge \left|\left(\sum_{1\le i\le k:f^m(x_i)=x_j}a_i-3^{k+j}\right) \right|-\left|\left(\sum_{1\le i\le N:f^m(x_i)=x_j}a_i-3^{k+j}\right)\right|\\
&>2-1=1.
\end{align*}


We can now conclude from $(2)$ that for every $1\le j\le k$, there exists $i\in ]k,N]$ such that $f^m(x_i)=x_j$ and $d(x_i,0)<\varepsilon$.\\

It remains to show that for every $1\le j\le k$, there exists $x\in M$ such that $d(x,0)<\varepsilon$  and $d(f^m(x),f^m(x_j))<\varepsilon$.  
Let $1\le j\le k$. 
If $d(f^m(x_j),0)<\varepsilon$ then it suffices to consider $x=0$. From now, we will thus assume that $d(f^m(x_j),0)\ge \varepsilon$ and we remark that it follows from $(1)$ that for every non-empty set $F\subset \{1,\dots,k\}$, we have
\begin{equation}\label{eq_coeff}
    \left|\sum_{i\in F}a_i\right|> 2.
\end{equation}

There are now two possibilities:
\begin{itemize}
\item If $f^m(x_j)\notin M_k$, then we deduce from $(4)$ that 
\[\left|\sum_{1\le i\le N:f^m(x_i)=f^m(x_j)}a_i\right| d(f^m(x_j),0)<\varepsilon\]
and since we have assumed that $d(f^m(x_j),0)\ge \varepsilon$. We deduce that 
\[\left|\sum_{1\le i\le N:f^m(x_i)=f^m(x_j)}a_i\right|<1.\]
Combined with \eqref{eq_coeff} for $F=\{1\le i\le k:f^m(x_i)=f^m(x_j)  \}\ni j$, it follows that  
\[\left|\sum_{i\in ]k,N]:f^m(x_i)=f^m(x_j)}a_i\right|\ge 1.\]
Therefore, by $(2)$, there exists $i\in ]k,N]$ such that $f^m(x_i)=f^m(x_j)$ and $d(x_i,0)<\varepsilon$. Thus we can consider $x=x_i$ in this case.

\item If $f^m(x_j)\in M_k$, then $f^m(x_j)=x_J$ for some $1\le J\le k$ and we use $(3)$ to get that 
\[\left|\sum_{1\leq i\leq N:f^m(x_i)=f^m(x_j)}a_i-3^{k+J}\right|<1.\]
We can then deduce from $(1)$ that
\[\left|\sum_{i\in ]k,N]:f^m(x_i)=f^m(x_j)}a_i\right|\ge 1\]
thus by $(2)$, there exists $i\in ]k,N]$ such that $f^m(x_i)=f^m(x_j)$ and $d(x_i,0)<\varepsilon$. Then, we can consider $x=x_i$ in this case.
\end{itemize}

\end{proof}

\end{document}
